%% 第一章 绪论
\chapter{绪论}

\section{研究背景}

\subsection{低空经济发展与无人机规模化应用}

过去十年间，微机电系统、飞控算法、电池与无线通信等技术的成熟，使多旋翼无人机从专业航拍设备迅速演变为大众消费品。以 DJI、道通、零度等厂商为代表的企业将整机价格下探至数千元量级，起飞重量低于 250\,g 的微型无人机甚至可以免注册购买。与之相伴的是低空空域活动密度的快速上升：航拍测绘、电力巡检、农林植保、快递物流、应急救援、灯光表演等应用场景不断扩展，城市低空正在形成一个新的经济形态，即通常所说的"低空经济"。

低空经济的繁荣带来了空域治理的新问题。一方面，无人机的合法运行需要飞行计划报备、空域申请、实名登记与实时监管；另一方面，无人机的获取门槛极低、操控简单、体积小巧，使得"黑飞"（未经许可的飞行）、"乱飞"（违反禁飞区或高度限制的飞行）以及恶意入侵难以通过传统手段有效遏制。近年来，国内外多次出现无人机闯入机场净空区导致航班大面积延误、无人机近距离拍摄军事设施或敏感区域、无人机携带危险载荷飞越人群等事件。这些事件表明，低空安全已经从理论风险转变为现实威胁，反无人机（Counter-UAV，C-UAV）技术因此成为学术界、产业界与监管部门共同关注的热点。

\subsection{反无人机系统的典型环节与技术路线}

一个完整的反无人机系统通常可以划分为"探测—识别—跟踪—定位—处置"五个环节\cite{radar_cuav,fusion_cuav}。探测环节负责发现空域中是否存在可疑目标；识别环节负责判定目标类型（无人机、鸟类、风筝、飞机等）；跟踪环节负责在时间序列上维持对同一目标的连续观测；定位环节负责给出目标相对观测平台或地理坐标系的位置；处置环节则根据威胁等级采取告警、干扰、诱骗或拦截等措施。其中，探测、跟踪与定位构成了整个系统的信息基础，其性能直接决定了后续处置是否可行。

按照所使用的物理媒介，现有探测手段可以分为雷达、射频、声学、光学以及多传感器融合五类，表~\ref{tab:sensing_compare} 对它们进行了对比。

\begin{table}[htbp]
\centering
\caption{典型反无人机探测手段对比}
\label{tab:sensing_compare}
\small
\begin{tabularx}{\textwidth}{P{2.0cm}P{2.2cm}P{2.5cm}P{2.8cm}X}
\toprule
\textbf{手段} & \textbf{作用距离} & \textbf{优势} & \textbf{局限} & \textbf{适用场景} \\
\midrule
雷达 & 数百米至数公里 & 全天候、测距测速直接 & 小型目标雷达散射截面积小、设备昂贵、需频率许可 & 机场、边境、要地远程预警 \\
射频侦测 & 数百米至数公里 & 被动工作、可识别型号 & 目标静默时失效、无法定位非合作目标 & 城市频谱监测 \\
声学探测 & 数十米至数百米 & 成本低、可被动布设 & 易受环境噪声干扰、作用距离短 & 楼宇近程补盲 \\
光学视觉 & 数十米至数百米 & 成本低、部署灵活、可直观确认 & 受光照与天气影响、小目标检出难、遮挡敏感 & 近程确认、跟踪与取证 \\
多传感器融合 & 视配置而定 & 互补性强、可靠性高 & 系统复杂、时空配准与数据融合难度大 & 高等级防护系统 \\
\bottomrule
\end{tabularx}
\end{table}

由表可见，雷达适合远距离预警，但对低慢小目标的探测能力有限，且设备成本与部署门槛较高\cite{radar_cuav}；射频侦测属于被动手段，一旦目标进入静默状态便完全失效\cite{rf_drone}；声学探测成本低廉，但作用距离通常只有百米量级，且在城市噪声环境下虚警率较高\cite{acoustic_drone}；多传感器融合虽然性能最好，却需要解决异构传感器的时空配准、数据关联与决策融合等一系列工程问题\cite{fusion_cuav,rf_remoteid}。相比之下，光学视觉方案只需要一台普通可见光相机与一台计算设备，成本低、部署灵活，并且能够提供目标的外观证据，便于操作员人工复核，因而成为近程反无人机系统中不可替代的组成部分\cite{optical_cuav}。

\subsection{光学视觉方案面临的核心困难}

尽管优势明显，光学视觉方案在真实反无人机场景中仍面临两方面突出困难。

第一，\textbf{小目标检测难}。反无人机系统追求"早发现"，而越早发现意味着目标越远、成像越小。在 1280$\times$720 分辨率的图像中，100\,m 外的微型无人机往往只占十几个像素，即使在大倍率长焦镜头下也可能不足 50 像素。通用目标检测模型在 COCO 等以中大型目标为主的数据集上训练，其特征金字塔的最深层对应 32 倍下采样，微小目标在深层特征图上仅剩不到一个像素，响应极弱。Det-Fly 等无人机检测数据集的统计也表明，目标框面积中位数不足画面面积的 0.06\%，99\% 以上的目标框面积小于画面的 1\%。如何在保持实时性的前提下提升小目标召回率，是检测环节的首要问题。

第二，\textbf{遮挡与抖动导致跟踪中断}。真实低空空域中充满树木、建筑、电线、广告牌等遮挡物，目标在飞行过程中频繁被短时遮挡。若系统以"每帧检测成功"为输出前提，则遮挡发生时输出会瞬间消失，航迹断裂、告警闪烁，操作员无法判断目标是暂时被遮挡还是已经飞离。此外，检测框本身存在像素级抖动，直接由框尺寸反算距离会产生米级跳变。如何在检测缺失时维持航迹、在检测存在时平滑输出，是跟踪与定位环节的核心问题。

除上述两点之外，工程落地还面临第三个困难：\textbf{仿真与真实域之间的差异}。出于成本与空域限制，算法开发通常先在仿真环境中完成闭环验证，再迁移到真实设备。然而仿真渲染的光照、纹理、边缘与真实相机存在显著差异，在真实数据上训练的模型直接用于仿真画面时可能出现严重的性能退化。本文在实验初期即观察到了这一现象，并将其作为一项专门的研究内容。

\section{研究意义}

\subsection{理论意义}

本文的研究意义首先体现在对小目标检测机理的补充上。现有工作大多从"增加检测层""引入注意力""改进损失函数"等角度提升指标，但对改进为何有效、在何种条件下失效缺乏统一解释。本文通过四组消融实验，将 P2 检测头的作用归结为"保留浅层空间分辨率"，将 EMA 的作用归结为"在通道与空间两个维度重加权"，并通过 +EMA 单独失效、+P2+EMA 联合有效的现象，说明二者存在"看清"与"认准"的互补关系。这一解释对小目标检测的网络设计具有一般性参考价值。

其次，本文对仿真域检测死区的机理分析具有方法论意义。面对"模型在仿真中失效"这一现象，容易产生的第一反应是"视角不合适"或"渲染质量差"。本文通过真值几何审计与距离分层统计，定量排除了视角假设，将原因定位为目标尺度分布偏移，并进一步用仿真域微调实验验证了该解释。这种"现象—假设—证伪—验证"的分析路径，比单纯报告指标更有助于理解域偏移的本质。

最后，本文将三态状态机、单目测距与自适应滤波整合为一条可运行的完整链路，并在统一的时间基准下评估检测、跟踪与定位的耦合性能，为反无人机视觉系统提供了一套可复现的工程范式。

\subsection{工程意义}

从工程角度看，本文系统具有以下价值。其一，硬件门槛低。系统仅依赖单目可见光相机与一台带 GPU 的普通计算机，不需要激光雷达、毫米波雷达或高精度惯导，便于在园区、楼宇、临时活动现场快速部署。其二，链路完整。系统不仅输出检测框，还输出目标距离、方位角、俯仰角、速度、威胁等级以及地理坐标，可直接对接上级指挥系统。其三，可测试性强。本文设计了 14 个自动化用例与统一的 CSV 记录格式，能够在无人值守条件下完成回归测试，为后续迭代提供量化基线。其四，可复现性强。仿真环境、数据采集脚本、自动标注流程与训练配置均已固化，研究者和工程师可以在相同条件下复现本文的全部实验。

\section{国内外研究现状}

\subsection{通用目标检测算法的发展}

目标检测算法的发展大致经历了三个阶段。第一阶段以手工特征与分类器组合为主，代表性方法是 HOG+SVM 与可变形部件模型（DPM），其检测精度依赖人工设计的特征，泛化能力有限。第二阶段以深度卷积网络为特征提取器，形成了两阶段与单阶段两条技术路线。两阶段方法以 Faster R-CNN\cite{faster_rcnn} 为代表，先由区域建议网络（RPN）生成候选框，再对候选框进行分类与回归，精度较高但速度较慢；单阶段方法以 SSD\cite{ssd} 与 YOLO 系列\cite{yolov1,yolov5,yolov8} 为代表，直接在特征图上预测类别与边框，速度优势明显。第三阶段以 Transformer 架构为代表，RT-DETR\cite{rt_detr} 通过编码器—解码器结构与集合预测损失摆脱了锚框与非极大值抑制，在精度上接近甚至超过同期 YOLO，但计算量与部署复杂度更高。

在工程实践中，YOLO 系列因其"精度—速度—易用性"的平衡而成为主流选择。YOLOv5\cite{yolov5} 确立了以 CSPDarknet 为骨干、PANet 为颈部、解耦头为头部的结构；YOLOv8\cite{yolov8} 在此基础上引入 C2f 模块与更灵活的任务对齐分配策略，并提供 n/s/m/l/x 五种规模。本文选择 YOLOv8n 作为基础模型，主要基于三点考虑：参数量约 3\,M，便于与桌面应用集成；在 RTX 4060 级别 GPU 上批量推理可达 100\,fps 以上；Ultralytics 提供了成熟的训练、验证与导出接口，便于复现实验。

\subsection{无人机小目标检测研究现状}

针对无人机小目标检测，现有研究大致可以分为四类。

第一类是\textbf{增加高分辨率检测层}。标准 YOLO 的最浅检测层为 P3（stride=8），对于十几个像素的目标仍显不足。SOD-YOLOv8\cite{sod_yolov8} 在 YOLOv8 中增加更高分辨率的检测分支，并配合特征融合，提升了交通场景中小目标的召回率。本文将这一思想应用于无人机场景，增加 stride=4 的 P2 检测头，使 40\,px 的目标在特征图上仍占据 10 个网格单元，而不是 5 个。

第二类是\textbf{引入注意力机制}。CBAM\cite{cbam} 通过通道注意力与空间注意力串联，增强目标区域响应；EMA\cite{ema} 则采用分组并行与跨空间学习，在较低计算开销下同时建模通道与空间关系。YOLOv8-MPEB\cite{yolov8_mpeb} 将骨干替换为 MobileNetV3 并引入高效多尺度注意力，实现轻量化无人机检测。MSDF-Net\cite{msdf_net} 结合可变形卷积与高分辨率特征增强，在林业病害检测中验证了跨版本特征增强的有效性。

第三类是\textbf{改进损失函数与标签分配}。小目标对定位误差极为敏感，IoU 系列损失在目标框很小时梯度不稳定。SIoU\cite{seiou} 引入角度代价，使回归方向更明确；Wise-IoU 通过动态聚焦机制降低低质量样本的梯度权重。本文在训练中沿用 YOLOv8 默认的 TaskAlignedAssigner 与 CIoU 损失，未对损失函数做额外改动，以保持与基线模型的可比性。

第四类是\textbf{数据集与评测协议}。Det-Fly\cite{detfly} 是目前公开的较大规模无人机检测数据集之一，包含多种背景、多种机型与多种飞行状态。UAVDT\cite{uavdt} 则面向城市交通场景中的无人机跟踪。Kurmashev 等人\cite{yolov12_adbc} 系统比较了不同 YOLO 版本在无人机检测任务上的表现，指出模型规模、输入分辨率与数据分布对结果的影响远大于网络结构的细微差异。这一结论提醒我们：在报告指标时必须明确数据划分与评测口径，否则容易得出误导性结论。

\subsection{反无人机系统研究现状}

完整的反无人机系统研究主要集中在三个方向。

其一是\textbf{多传感器融合}。文献\cite{rf_remoteid} 提出融合射频、光学与雷达的多模态框架，利用 Remote ID 广播信号进行身份识别，利用视觉进行确认，利用雷达进行测距；文献\cite{fusion_cuav} 对现有反无人机方法与基准进行了系统综述，指出单一传感器在复杂环境下均存在盲区，融合是提升可靠性的必然选择。

其二是\textbf{射频与声学探测}。文献\cite{rf_drone} 利用无人机图传与遥控链路的射频特征实现检测与跟踪，能够在目标静默前提供较远距离的预警；文献\cite{acoustic_drone} 利用 MEMS 麦克风阵列与机器学习实现被动声学定位，成本极低但作用距离受限。

其三是\textbf{视觉检测与跟踪}。文献\cite{optical_cuav} 对比了 YOLO 与 Faster R-CNN 在监控视频中的微型无人机检测性能，发现单阶段模型在速度上具有明显优势，而两阶段模型在极小目标上略优。总体来看，现有工作多聚焦于单一算法环节的指标提升，面向完整链路、可实际运行、具备可视化界面的系统级开源实现相对缺乏，这正是本文的切入点。

\subsection{单目定位与目标跟踪研究现状}

在定位方面，单目视觉测距通常有三类方法。第一类基于已知目标尺寸，利用相似三角形关系由框高或框宽反算距离，实现简单但对目标尺寸先验依赖强；第二类基于地面平面假设，利用相机高度与像平面几何关系解算目标落点，适用于地面目标；第三类基于多帧运动，通过特征点跟踪与三角化恢复深度，计算量大且对纹理敏感。本文的目标是空中无人机，不具备地面平面约束，且外观尺寸相对稳定，因此选择第一类方法，并通过仿真域标定获得等效尺寸。

在跟踪方面，SORT\cite{sort} 用卡尔曼滤波预测目标位置、用匈牙利算法关联检测框，奠定了多目标跟踪的基本框架；DeepSORT\cite{deepsort} 引入外观特征以降低身份切换。本文面向单目标跟踪场景，不需要外观重识别，但借鉴了"预测—关联—更新"的思想，并将关联失败时的处理显式建模为 COAST 状态。卡尔曼滤波\cite{kalman} 是最优线性估计的经典方法，但其量测噪声协方差需要预先设定。Sage-Husa 自适应滤波\cite{sagehusa} 通过残差在线估计噪声统计特性，能够在量测质量变化时自动调整滤波增益，本文采用该思路构建自适应卡尔曼滤波器（AKF）。

\subsection{仿真验证与域适应研究现状}

受空域管制与成本限制，无人机算法的早期验证通常依赖仿真。PX4\cite{px4} 是应用最广泛的开源飞控固件，支持 SITL（Software In The Loop）模式；jMAVSim\cite{jmavsim} 是其自带的轻量级三维仿真器，基于 Java3D 构建，能够以较低的资源开销模拟多旋翼动力学与相机视图。Gazebo 与 AirSim 提供更逼真的渲染，但对显卡与计算资源要求更高。本文选择 jMAVSim，主要原因是其启动快、可定制性强，便于通过 Java 源码改造注入目标机、相机与真值广播。

仿真到真实的域差异（sim-to-real gap）是近年来的研究热点。常见做法包括风格迁移、域随机化与目标域微调。风格迁移通过生成对抗网络将仿真图像转换为真实风格，但可能引入伪影；域随机化在训练时随机改变光照、纹理与噪声，提升模型鲁棒性，但需要大量训练资源；目标域微调则直接用少量目标域数据调整模型参数，简单有效。本文采用目标域微调路线，并针对仿真环境设计了零人工标注的采集流程，以较低成本消除了检测死区。

\section{现有方案的不足与本文切入点}

综合上述分析，现有方案存在以下不足。

（1）\textbf{小目标改进缺乏机理层面的解释}。多数工作只报告 mAP 提升，未说明改进在何种尺度、何种背景下起作用，也未讨论单独使用某一模块时可能出现的性能下降。本文通过消融实验与特征尺度分析，给出 P2 与 EMA 互补的解释。

（2）\textbf{仿真验证中的性能退化被忽视}。不少工作在仿真环境中直接使用真实数据训练的模型，遇到性能下降时归因于"渲染不真实"，缺少定量分析。本文通过逐帧真值审计定位死区成因，并用微调实验闭环验证。

（3）\textbf{跟踪与定位耦合不足}。检测、跟踪、测距往往被分别评估，缺少在统一时间基准下的端到端指标。本文在仿真闭环中以 50\,Hz 真值为基准，同时评估检出率、测距误差、角度误差与预测误差。

（4）\textbf{系统级实现与可复现性不足}。许多研究停留在脚本与 notebooks 层面，缺少可视化界面、参数管理、日志记录与自动化测试。本文构建了完整的桌面系统与测试框架，并将全部配置固化。

基于以上分析，本文的切入点可以概括为：\textbf{以单目可见光相机为传感器，以改进的 YOLOv8n 为检测前端，以三态状态机为跟踪骨架，以自适应卡尔曼滤波为定位核心，以仿真闭环为验证平台，构建一套端到端可运行、可测试、可复现的反无人机识别与定位系统}。

\section{本文主要研究内容}

本文的主要研究内容包括以下五个方面。

（1）\textbf{小目标检测模型改进与消融分析}。在 YOLOv8n 上增加 P2 检测头与 EMA 注意力，构建四种模型配置，在 Det-Fly 数据集上按序列划分训练/验证集，完成消融实验，并从特征分辨率与注意力重加权两个角度解释结果。

（2）\textbf{数据集构建与数据工程}。完成 Det-Fly 数据格式转换、序列级划分与统计分析；采集自建 AprilTag 标定数据用于相机内参标定与定位精度评估；构建 450 帧仿真自动标注数据集，用于仿真域微调。

（3）\textbf{三态跟踪状态机设计}。针对误检、跳变与遮挡三个问题，设计 SEARCH/LOCKED/COAST 状态机，给出双门限、空间门限与超时参数的选择依据，并通过遮挡用例验证 COAST 的航迹保持效果。

（4）\textbf{单目定位与自适应滤波}。建立针孔相机模型，推导方位角、俯仰角与距离的解算公式，分析测距误差随距离与框高误差的传播规律；设计六维状态自适应卡尔曼滤波器，给出预测步与更新步的完整推导，并通过仿真与实拍数据评估平滑与预测性能。

（5）\textbf{仿真闭环与系统实现}。改造 PX4+jMAVSim 闭环，注入脚本目标机、护航相机、真值广播与离屏帧流；分析并消除检测死区；基于 PySide6 实现三页签监视大屏；设计 14 个自动化用例完成回归测试；进一步实现目标机主动飞行、相机被动尾随的编队巡游与加减速行为。

\section{论文组织结构}

本文共分十章，各章内容安排如下。

第一章为绪论。介绍研究背景、反无人机系统的技术路线、光学视觉方案的核心困难，综述国内外研究现状，分析现有方案的不足，并给出本文的研究内容与组织结构。

第二章为相关理论与关键技术。介绍卷积神经网络与目标检测的基本原理、YOLOv8 的网络结构、注意力机制、卡尔曼滤波、针孔相机模型、PX4 与 jMAVSim 仿真框架、PySide6 与网络通信等本文所依赖的理论与技术基础。

第三章为数据集构建与数据工程。介绍 Det-Fly 数据集的来源、格式、统计特性与序列级划分方法；介绍自建 AprilTag 数据集的采集与标定流程；介绍仿真自动标注数据集的构建方法；讨论数据增强与评测口径。

第四章为检测模型改进与消融实验。详细阐述 P2 检测头与 EMA 注意力的设计，给出网络结构、计算量与感受野分析；介绍训练配置与损失函数；给出四组消融实验结果，并从尺度分布、注意力响应与参数效率三个角度解释结果；最后讨论模型的有效工作范围。

第五章为三态跟踪状态机。分析直接使用检测输出导致的误检、跳变与遮挡问题；给出状态定义、转移条件与参数选择；详细描述 SEARCH、LOCKED、COAST 三个状态的处理流程；给出伪代码与复杂度分析；通过遮挡用例验证航迹保持效果。

第六章为单目定位与自适应卡尔曼滤波。建立相机坐标系与像素坐标系的映射关系；推导方位角、俯仰角与距离的解算公式；分析测距误差传播；设计六维状态自适应卡尔曼滤波器，推导预测与更新过程；给出角度解算的两种实现及其对比；最后通过仿真与实拍数据评估定位精度。

第七章为仿真闭环系统与域微调。介绍 PX4 SITL 与 jMAVSim 的架构；详细说明本文对 jMAVSim 的改造，包括脚本目标机、护航相机、真值广播、离屏帧流与场景构建；描述检测死区的发现过程与机理分析；给出自动标注数据采集、距离分层采样与微调配置；展示微调前后的对比结果；介绍编队巡游与尾随相机的实现。

第八章为系统设计与实现。给出需求分析、总体架构与双进程通信协议；介绍算法端的模块划分、状态机实现与数据记录；介绍显示端的界面布局、地图投影、光电画面、仪表与告警；介绍参数管理、回放与自动化测试框架；讨论部署与运行维护。

第九章为实验与结果分析。给出实验环境与评测口径；分别报告检测模型性能、仿真闭环定位精度、实拍视频定位精度、退化鲁棒性与 14 用例回归测试结果；进行横向对比与消融分析；讨论系统能力边界与失败案例。

第十章为总结与展望。总结全文工作与创新点，分析存在的不足，展望后续研究方向。
